\documentclass[a4paper,12pt]{ltjsarticle}

% --------------------
% 必須最低限
% --------------------
\usepackage{amsmath,amssymb}
\usepackage{geometry}
\usepackage{enumitem}
\usepackage{xcolor}
\usepackage{luatexja-fontspec}
\usepackage{hyperref}

% --------------------
% レイアウト（書き足しやすさ重視）
% --------------------
\geometry{
  top=22mm,
  bottom=22mm,
  left=30mm,
  right=30mm
}

\renewcommand{\baselinestretch}{1.3}
\setlength{\parskip}{0.7em}
% \setlength{\parindent}{0pt} % ノート感を出すため字下げなし

% --------------------
% フォント（無難・疲れにくい）
% --------------------
\setmainjfont{IPAexMincho}
\setmainfont{Times New Roman}

% --------------------
% 箇条書き（軽量）
% --------------------
\setlist[itemize]{
  label=--,
  leftmargin=1.5em,
  itemsep=0.2em
}

% --------------------
% メモ用コマンド
% --------------------
\newcommand{\memo}[1]{%
  \par\noindent
  {\color{gray}\small [#1]}
}

\newcommand{\todo}[1]{%
  \par\noindent
  {\color{blue}\bfseries TODO: #1}
}

% --------------------
% 見出しを弱くする
% --------------------
% \usepackage{titlesec}
% \titleformat{\section}
%   {\normalfont\bfseries}
%   {}
%   {0pt}
%   {}
% \titlespacing*{\section}{0pt}{1.5ex}{0.8ex}

% --------------------
% ハイパーリンク（控えめ）
% --------------------
\hypersetup{
  colorlinks=true,
  linkcolor=black,
  urlcolor=blue
}

%====================
% 強調用コマンド
%====================
\newcommand{\important}[1]{\textcolor{blue}{\textbf{#1}}}

\begin{document}

\begin{center}
  {\huge 考えのまとめ}\\[1em]
  {\large \today}\\
  {\large 千葉工業大学大学院 学習工学研究室}\\
  {\large 内山 裕太}
\end{center}

\section{悩みの背景}

東本研究室では，東本先生なしでも十分に回る研究室運営が重要とされている．これは，指導教員である東本先生が「学生一人ひとりが成長できる場であってほしい」という考えを持っているからであり，学生は主体的に考えて行動することが求められる．個々人が自身の成長を考えた上で，研究室の中でどのように活動するべきかを考える必要がある．\\

また，研究室内には素晴らしい後輩がたくさんいる．自分の弱点を成長させようと頑張る人や，自身の成長のために率先して動く人などそれぞれである．しかし，個人の成長を重要視して動くこと自体は十分に納得できる一方で，研究室という「集団」として知見を共有していくという視点では課題がある気がしている．そもそもとして，\important{自分は博士課程に進学するわけではないため，研究室の将来を強く背負う責任はないが，だからと言って研究室としてのあり方に無関心で良いという考え方は倫理的にまずいと思うし，自分がいる期間に共有できる考え方は共有しあいたいと思っている．}\\

一方で，自分は塾講師として約5年働いている．教育を経由して，一人ひとりの生徒の将来に貢献したいという気持ちがあり，そのまま同じ企業に内定をもらっている．校舎には校舎長がおり，基本的には校舎長の方針に沿って動くことが前提となる．
\newpage

\section{問題点}

では，何が問題なのか？それは

\important{環境ごとに求められる役割や価値観が異なることと，それに自分が柔軟に対応するための考え方を持ち合わせていないこと}\\

研究室では，主体的に考え，学生自らが課題を見つけて行動する必要がある．しかし，就職先ではトップの考え方を踏まえた上で，組織としての一貫性を保ちつつ働く必要がある．だからこそ，研究室活動の中で育ってきたような「もっとこうした方が良いのでは？」という考え方自体が，そのまま出しにくい場面があると思う．\\

さらに面倒なことに，自分は研究室にあと1年くらいしか残らない橘にも関わらず，集団としてのあり方に違和感を持ってしまっている．これ自体は悪いことだと思っていないが，\important{「自分はそこまで責任を持つ必要がない」}という考えと\important{「違和感や問題に気づいている以上，何もしないのは違うのでは？」}という考えの間で彷徨い続けている．\\

問題が\important{単なる環境の違い}ではなく，\important{自分はどこまで主体的に関わるべきなのか，どこまで責任（責任感）を持つべきなのか}が定まっていないことである．
\newpage

\section{現状の良いところとは？}

自分なりに，自分の考え方の良いところを考えてみた

\begin{itemize}
  \item 主体的に考えを持っていること
  \item 集団としてみれる視野の広さ
  \item 表現し難い悩みを外化しようとしていること
  \item 自身がどうあるべきかを客観的に見ようとしていること
\end{itemize}

一般的に考えれば，人から言われたことを忠実にこなしていれば良い．ただ，研究活動を通して「本当にこれで良いのか」「本来はどうあるべきなのか」を考えることができている点は良いところだと考えている．また，自分の倫理的な違和感をベースに，集団としてのあり方や後輩の行動が本当に全体の成長につながるのかという視点で客観的に見れている点も良い．
\newpage

\section{何を直し，調整するべきか？}

\subsection{主体性と役割の意識を分けて考えられるように}

今の自分は\important{「自分で考えること」}が強く出過ぎていることから\important{「自分の考えを通すこと」}に変わってしまっているところがある．逆に，このような考え方をできているからこそ\important{その場ごとに何が求められていて，自分がどうあるべきか，何をするべきなのか}をしっかり考えて動く必要がある．\\

研究室では，後輩一人ひとりがどのように考えて行動しているのかをしっかりと理解し，研究室全体としてみた時に本当にその考え方で大丈夫なのかを共有する必要がある．自分1人で問題を抱え込むのではなく，集団全体に問題提起をしていく振る舞い方をする必要がある（すでにやってる？）．

\subsection{責任ではなく責任感}

研究室として大丈夫？などの考え方を持つこと自体は大切だが，それが改善できない責任が自分にあるわけではない．確かに，将来的なことを考えると自分はそこまで関与しなくても良いかもしれない．ただ，問題として提起して，集団としてその問題に向き合うことは重要．仮に解決にならなかったとしても，それは自分の責任じゃない．だからこそ，疑問に思ったことなどを全体に共有する考え方を持った方が良い．

\subsection{「正しさ」を一つに固定しない}

研究室，就職先での考え方は矛盾する．ただ，これはあくまでも大枠の話であり，実際は場面によって考え方や動き化やが異なるだけ．無理やり統一するのではなく，それぞれの場面での経験を通して，臨機応変により良い結果へと結びつけられるような振る舞いを学ぶべき．
\newpage

\section{今後をどうするか}

\subsection{就職先}

本日（3/11）に，校舎長に自分なりの考えを伝えました．\important{「方針を決めるのは校舎長だからこそ，意図を理解して，最大限の価値を見出せるように動きたい．その辺を整理したいから，会議の前に少し話す時間をほしい」}という内容です．3/24の会議の前に，しっかり1年の方針を決める時間をもらえました．\\

間違っていると思うことは多々ありますが，否定的な目線で入るのではなく，提案として伝えられるような姿勢を身につけたいです．\important{「自分の考えや提案を殺す」のではなく，「組織の状態に合わせて意見の出し方を調整する」ということを意識します．}

\subsection{研究室}

研究室では，ラスト1年という立場を踏まえつつ，無関心にはならずに\important{自分の関われる範囲でより良い姿勢を示せるようにすることを意識します．}\\

何か違和感を感じるなら，全てを一気に全面的に変えようとするのではなく，\important{集団としてどうした方が良いのか}という自分の意見をしっかり他者に伝える．この考え方や動き方は今後も続けていこうと思いますが，院生になる4人を踏まえて，今の組織が組織としてどう思っているのかを一度聞く必要がありそうです．自分は問題だと思っていても，そこまで問題に思っていない人だらけかもしれませんから（先生を含めてM以上はわかってくれていると思ってます）．

\subsection{内山として}

今までに色々な面で自分の考え方を共有してきたつもりでした．プレゼミで寝てしまう人がいることに対しては「学会同様に議論スレッドを立てる」，1人が率先して動きすぎてしまっていることに対しては「院生ディスカッションのある意味について，自分なりの意見を提示」などです．\important{後者に対しては，あまり返信がなくて悲しかったですが，今後も自分なりの考えは共有します．}\\

研究室と就職先では，役割や考え方が異なります．主体性を持ちながらも，場に応じた役割調整をできるようになることが課題です．また，問題に気づいた時に，全てを背負うのではなく，\important{「自分の立場であればどこまで関わることができるのか」，「まずは何をするべきなのか」}を考えて行動に移す必要があります．目標は\important{「自分の価値観を保った上で，環境に応じた適切な関わり方を切り替えられる人間になる」}です．\\

あまり気負いしすぎずに，自分で決めた時間に研究室に行って，その日にやるべきことをやって，満足できる状態を作ること．全てを自分で考えようとするのではなく，自分の考え方等を共有しながらも，責任を感じすぎて研究室に行くのが嫌になったら本末転倒なので，自分に責任はないという考えをしっかり持って，責任感のある1年間にできたらと思います．
\newpage

\section{さいごに}

東本先生，前田先輩，ピエレット君にはいつもご迷惑と心配ばかりかけてしまい申し訳ないです．ここ2週間くらい考えて，自分はどうあるべきなのかをずっと悩んでいました．東本先生からのメッセージを読んで，少なくとも自分の考え方に大きな問題があるわけではないと思えました．ただ，まだまだ自分のマネジメント力は足りないところがあり，成長すべき点が山積みだと感じました．\\

人聞き悪い言い方になりますが，東本先生の「学生主体の研究室運営」は，一見すると「指導教員が楽をしているだけ」とみられるかもしれませんが，非常に難しい運営方針だと痛感しました．本来，自分が学生個々人をマネジメントすれば済むところを，学生が成長できるレールを敷ける環境を構築しながら他人に行わせるという非常に難しいことを行なっていると思います．本当にいつもありがとうございます．\\

3人には今後もご心配やご迷惑をおかけするかもしれません．1つ1つの経験や考え方の対立，吸収が自分のことを成長させてくれています．院ゼミなどを通して，今後も自分の悩み等については共有していけるように頑張ります．\\

長くなりましたが，以上が自分の考えのまとめになります．

\rightline{\large 内山}

\end{document}
